\documentclass[12pt, a4paper]{article}
\usepackage[utf8]{inputenc}
\usepackage[T1]{fontenc}
\usepackage[spanish,es-noshorthands]{babel}
\usepackage[margin=2.5cm]{geometry}
\usepackage{amsmath}
\usepackage{url}
\usepackage{booktabs}
\usepackage{xcolor}
\usepackage{tcolorbox}
\usepackage{listings}
\usepackage[colorlinks=true, linkcolor=blue!60!black, urlcolor=blue!60!black]{hyperref}

\definecolor{azulTitulo}{RGB}{20,60,120}

\lstset{
  basicstyle=\footnotesize\ttfamily,
  breaklines=true,
  breakatwhitespace=true,
  frame=single,
  framesep=4pt,
  backgroundcolor=\color{gray!8},
  literate={á}{{\'a}}1 {é}{{\'e}}1 {í}{{\'i}}1 {ó}{{\'o}}1 {ú}{{\'u}}1
           {ñ}{{\~n}}1 {Á}{{\'A}}1 {É}{{\'E}}1 {Í}{{\'I}}1 {Ó}{{\'O}}1 {Ú}{{\'U}}1 {Ñ}{{\~N}}1
}

\title{\textbf{Pipeline de Captura y Publicación de Datasets para Entrenamiento de LLMs}\\[6pt]
{\large GIDEAL — Workspace ELECTRONICA}}
\author{Prof. César Rodríguez}
\date{6 de agosto de 2026}

\begin{document}

\maketitle

\begin{tcolorbox}[colback=blue!5, colframe=azulTitulo, title=Resumen ejecutivo]
Se desarrolló e integró una \textbf{cadena completa de valor de datos} dentro del
workspace ELECTRONICA: la captura pasiva de muestras de entrenamiento generadas por
los flujos existentes (RAG, análisis de imágenes y conversión imagen a KiCad), su
curado y particionado, su empaquetado en formato Hugging Face con metadata de
licencia, y su publicación en el Hugging Face Hub. Todo siguiendo la arquitectura
de tres capas del proyecto: directivas, orquestación y ejecución.
\end{tcolorbox}

\section{Arquitectura de 3 capas}

\begin{center}
\begin{tabular}{lll}
\toprule
\textbf{Capa} & \textbf{Dónde} & \textbf{Rol} \\
\midrule
Directivas & \path{directives/*.yaml} & SOP: qué hacer \\
Orquestación & \path{flujo_*.py} (raíz) & Decisiones y validación \\
Ejecución & \path{execution/*.py} & Scripts deterministas: cómo \\
\bottomrule
\end{tabular}
\end{center}

La lógica probabilística (LLMs) se mantiene en los orquestadores; todo el
procesamiento de datos vive en scripts deterministas y reutilizables.

\section{Cadena de valor: 4 eslabones}

\begin{enumerate}
    \item \textbf{Captura pasiva} — registra muestras sin interrumpir los flujos
    \item \textbf{Curado} — filtra, valida, deduplica y particiona
    \item \textbf{Empaquetado} — genera paquetes Hugging Face con licencia
    \item \textbf{Publicación} — sube al Hugging Face Hub
\end{enumerate}

\subsection{1. Captura pasiva}

Tres orquestadores existentes fueron instrumentados para capturar sus salidas en
formato JSONL dentro de \path{datasets/} (excluido de git):

\begin{itemize}
    \item \path{rag_system.py} → \path{datasets/rag_conversaciones.jsonl} (ShareGPT)
    \item \path{flujo_imagen_a_kicad.py} → \path{datasets/eda_imagen_circuito.jsonl}
    \item \path{flujo_analizar_imagen.py} → \path{datasets/analisis_imagenes.jsonl}
\end{itemize}

La captura es 100\% pasiva: toda llamada va en \texttt{try/except}, sanitiza PII
(emails, cédulas, teléfonos $\rightarrow$ \texttt{[REDACTADO]}), descarta duplicados
por hash canónico y respeta la variable \texttt{DATA\_CAPTURE\_DISABLED} o el flag
\texttt{--no-capture-data}.

\subsection{2. Curado}

\path{execution/curar_datasets.py} lee los JSONL crudos y produce
\path{datasets/curated/}: valida el esquema, descarta registros con
\texttt{quality\_score} bajo el umbral (default 0.7), deduplica y particiona en
train/val/test (default 80-10-10, semilla fija 42). Emite \path{manifest.json} y
\path{resumen.md}.

\subsection{3. Empaquetado}

\path{execution/empaquetar_dataset.py} convierte los datasets curados en paquetes
formato Hugging Face en \path{datasets/paquetes/}: \path{dataset_info.json},
\path{README.md} (tarjeta con metadata), \path{LICENSE} y splits en Parquet o JSONL.
Soporta licencias \texttt{custom-proprietary} (GIDEAL-1.0), \texttt{apache-2.0},
\texttt{cc-by-4.0} y \texttt{mit}, además de empaquetado en \texttt{.tar.gz}
con el flag \texttt{--pack}.

\subsection{4. Publicación en Hugging Face Hub}

\path{execution/publicar_hf.py} lee \texttt{HF\_TOKEN} del entorno o de
\path{.env}, crea el repositorio de datasets (con \texttt{exist\_ok=True}) y sube
el contenido del paquete. Retorna la URL pública del repositorio. El flag
\texttt{--private} genera repositorios privados.

\section{Esquema de los datasets}

\begin{itemize}
    \item \textbf{RAG (ShareGPT):} \texttt{messages[]} (system/user/assistant) +
    \texttt{metadata.retrieved\_context[]}, \texttt{domain}, \texttt{model},
    \texttt{quality\_score}.
    \item \textbf{EDA imagen→circuito:} \texttt{images[]}, \texttt{instruction},
    \texttt{output} (components[] y connections[]) + \texttt{api\_backend},
    \texttt{model}, \texttt{tokens\_usados}.
    \item \textbf{Análisis multimodal:} \texttt{images[]}, \texttt{instruction},
    \texttt{output} (descripción, elementos detectados) + metadata de modelo.
\end{itemize}

Todo registro incluye \texttt{id}, \texttt{schema\_version} y \texttt{timestamp}.

\section{Comandos de uso}

\begin{lstlisting}[language=bash]
# Curado
python3 flujo_curar_dataset.py [--min-quality 0.7] [--split 80-10-10] [--dry-run] [--no-alert]

# Empaquetado
python3 flujo_empaquetar_dataset.py [--format parquet|jsonl] [--license <lic>] [--pack] [--no-alert]

# Publicación
python3 flujo_publicar_hf.py <dataset> [--repo <id>] [--private] [--dry-run] [--no-alert]
\end{lstlisting}

\section{Archivos creados y modificados}

\begin{itemize}
    \item \textbf{Creados:} \path{execution/data_capture.py},
    \path{execution/curar_datasets.py}, \path{execution/empaquetar_dataset.py},
    \path{execution/publicar_hf.py}, \path{flujo_curar_dataset.py},
    \path{flujo_empaquetar_dataset.py}, \path{flujo_publicar_hf.py} y cuatro
    directivas (\path{directives/data_capture.yaml}, \path{curar_dataset.yaml},
    \path{empaquetar_dataset.yaml}, \path{publicar_hf.yaml}).
    \item \textbf{Modificados:} \path{rag_system.py},
    \path{flujo_imagen_a_kicad.py}, \path{flujo_analizar_imagen.py},
    \path{.gitignore} (se excluye \path{datasets/}) y \path{AGENTS.md}.
\end{itemize}

\section{Prueba real publicada}

Se publicó el dataset de ejemplo \path{rag_conversaciones} en
\url{https://huggingface.co/datasets/cero2k6/gideal-rag-v1} con el token
recuperado de la caché de Hugging Face del usuario \texttt{cero2k6}. El
repositorio contiene \path{LICENSE}, \path{README.md}, \path{dataset_info.json}
y \path{train.parquet}, y queda como referencia de estructura. Se corrigió
durante el desarrollo que Hugging Face exige IDs de licencia válidos en el
frontmatter del README: las licencias propias se registran como
\texttt{other} + \texttt{license\_name}/\texttt{license\_link}.

\section{Comercialización (Web3 / Solana)}

El pipeline entrega los artefactos listos para la estrategia de monetización
definida en \path{Proyectos/SOLANA/docs/proyectos.md}. Hugging Face Hub es el
canal de \textbf{distribución}; la venta y verificación se canalizan en la red
Solana.

\begin{enumerate}
    \item \textbf{Dataset curado y empaquetado} — con Parquet, \path{LICENSE} y
    \path{README.md} con metadata (origen, licencia, versión, citación).
    \item \textbf{Almacenamiento descentralizado} — el paquete (o sus hashes) se
    sube a \textbf{Arweave} (permanente) o \textbf{IPFS} (con pinning en
    Pinata/Helius).
    \item \textbf{Tokenización} — un \textbf{Data NFT} en Solana representa el
    dataset: el contrato controla el acceso mediante \textbf{pago en USDC} (flujo
    x402) y distribuye royalties al wallet del autor.
    \item \textbf{Marketplaces} — Magic Eden, Tensor, Ocean Protocol o Kled AI
    según el tipo de dataset.
    \item \textbf{Monetización recurrente} — royalties on-chain (ej. 2--5\,\%)
    cada vez que un modelo entrenado con los datos se usa comercialmente.
\end{enumerate}

\subsection{Reglas de oro para datos en blockchain}

\begin{itemize}
    \item \textbf{Nunca subir PII on-chain}: la cadena es inmutable. En el
    contrato solo van hashes y metadatos; los datos crudos en almacenamiento con
    acceso controlado. El pipeline ya sanitiza PII en la captura.
    \item \textbf{Consentimiento}: si los datos incluyen interacciones con
    estudiantes, anonimizarlos u obtener consentimiento explícito
    (GDPR/CCPA).
    \item \textbf{Licencias claras en el contrato}: entrenamiento propio (sí),
    reventa (no), uso comercial (sí, con royalty), código abierto (negociable).
    La licencia GIDEAL-1.0 \path{LICENSE} del paquete es la base de estos
    términos.
    \item La \textbf{calidad} (metadato \texttt{quality\_score}) determina el
    precio: datasets curados valen más que crudos.
\end{itemize}

\section{Consideraciones y convenciones}

\begin{itemize}
    \item El directorio local \path{datasets/} \textbf{sombrea la librería}
    \texttt{datasets} de Hugging Face desde la raíz del workspace; los Parquet se
    cargan con \texttt{pyarrow} (ver README de los paquetes).
    \item \path{datasets/} y \path{.env} están excluidos de git
    (\path{.gitignore}).
    \item El token HF jamás se imprime ni registra; se lee de \path{.env} o del
    entorno.
    \item La licencia GIDEAL-1.0 no es de uso público libre: para datos comerciales
    se recomienda publicar con \texttt{--private}.
    \item Con hasta 5 muestras, todo el split va a train (val/test en 0); a partir
    de 6 muestras el particionado asigna 1 a val y 1 a test y se reparte
    automáticamente al acumular más datos.
\end{itemize}

\section{Trabajo futuro}

\begin{itemize}
    \item Acumular datos reales y publicar datasets con cientos de muestras.
    \item Automatizar el curado periódico (cron).
    \item Integrar la cadena con la estrategia Web3 de comercialización de datos
    (licencias on-chain, verificación).
    \item Añadir validación de calidad por heurísticas (longitud, consistencia del
    netlist, cobertura de imágenes).
\end{itemize}

\end{document}
